% AI-written, AI cross-reviewed; no human sign-off; not a claim that RH is proved
% Atlas unit T17. Local manuscript v3, item 5.27, heading on printed PDF p. 45.
% PositivityTheory.Pf1.HilbertSamuelMoriwaki.moriwakiBig_outright
% paper2/split_p012.tex:299–315. Verbatim excerpt; SHA-256 7e59d1ea740a3b8d05709705adf5c567a0fc00c6f921636f566dcb155f22e689.
% Original TeX uses the macros and cross-reference keys of the source manuscript.


% Start of atlas edition annotation. The following block is outside the original excerpt.
\paragraph{Status of named inputs in this reading.}
\noindent\nolinkurl{PositivityTheory.Pf1.HilbertSamuelMoriwaki.moriwakiBig_outright} --- in-repo constant (\path{Development/WeilBound/CoreStatements/AutomorphicKernels/MoriwakiBigness/MoriwakiBigTheorem/Outright.lean:9--11}); variant.\par
\noindent\nolinkurl{GlobalTheory.T.arith_intersection_system_of_poincare_lelong} --- in-repo constant (\path{Development/AutomorphicSpectrum/Proofs/MetricProofs/ArithIntersectionSystemOfPoincareLelongStatement.lean:11--12}); not compared.\par
\noindent\nolinkurl{PositivityTheory.Tls2.Pl.poincareLelong} --- in-repo constant (\path{Development/WeilBound/CoreStatements/TemperedRepresentationTheory/PoincareLelong/PoincareLelongInv.lean:13--17}); not compared.\par
% End of atlas edition annotation.

\begin{formalstatement}[Arithmetic intersection bounded by small-section volume]\label{sub:T17}

Take finite blocks of homogeneous coordinates $y$ in a field $\Omega$,
their vanishing ideal $P$, a birational presentation of arithmetic dimension
$D+1$, and an intersection system $I$ satisfying its defining axioms.
Let $A$ be a conjugation-invariant Hermitian line, semipositive relative
to some $P_0\subseteq P$, nef on the fibres in its multidegree and nef on
the generic fibre in the specified sense. Then
\[
 I_P(A^{D+1})\le
 \mathrm{arVolP}\bigl(\text{small sections for }P,y,A,\,
       \mathrm{chartLat}(y,\deg A),1,D\bigr).
\]

\end{formalstatement}\subnode{T17}

\noindent The relative semipositivity, birational chart lattice and arbitrary intersection system require comparison. \href{https://arxiv.org/pdf/math/0612269}{A. Moriwaki, Continuity of volumes on arithmetic varieties, Corollary 5.5}~\cite{r21-moriwaki}.
